% P05: REPLACES the TWO original figure blocks with this merged (A)/(B) figure.
% fig:Ext_passthru retired; both in-text refs repoint to fig:Ext10mm_52cm B.
\begin{figure}[htbp]
\begin{center}
\includegraphics[width=1\textwidth]{Ext10mm_52cm.eps}\par\smallskip
\includegraphics[width=0.52\textwidth]{KneeExt_10mm_passthru.png}
\end{center}
\caption{Measured versus expected isometric torque results using a $\diameter$\qty{10}{mm}, $l_{0}=\qty{51.8}{\cm}$ extensor BPA on the biomimetic knee (\textbf{A}), and the routing configuration that makes this test difficult to model (\textbf{B}). In \textbf{(A)}, measured results are blue dots, the original prediction is a dashed black line, the optimized prediction is a solid magenta line, and the hybrid calculation is shown in yellow dots. In \textbf{(B)}, the BPA passes between the femoral condyles close to the instantaneous center of rotation.}\label{fig:Ext10mm_52cm}
\end{figure}
